\documentclass[11pt]{article}
\usepackage[margin=1in]{geometry}
\usepackage{amsmath,amssymb,amsthm}
\usepackage{booktabs}
\usepackage{hyperref}

\newtheorem{theorem}{Theorem}
\newtheorem{lemma}[theorem]{Lemma}
\newtheorem{proposition}[theorem]{Proposition}
\newtheorem{corollary}[theorem]{Corollary}
\theoremstyle{definition}
\newtheorem{definition}[theorem]{Definition}
\theoremstyle{remark}
\newtheorem{remark}[theorem]{Remark}

\newcommand{\Lang}{\mathcal{L}}
\newcommand{\F}{\mathcal{F}}

\title{Disproof of Conjecture 00000009617:\\
the even shift has a 2-state Fischer cover but Hankel rank 3}
\author{jilint777}
\date{2026-10-04}

\begin{document}
\maketitle

\begin{abstract}
Conjecture 00000009617 asserts, among other things, that the Fischer
(``Fisher'') minimal state count of a sofic shift equals the rank of the
Hankel matrix of its language. We show that this clause is false. For the
even shift $X$ (binary sequences in which every run of $1$s between two $0$s
has even length) the Fischer cover has exactly $2$ states, and no
presentation of any kind has fewer. But the Hankel matrix
$H[u][v]=[uv\in\Lang(X)]$ has the $3\times 3$ minor
$\left(\begin{smallmatrix}1&1&1\\1&0&1\\0&1&1\end{smallmatrix}\right)$
with determinant $-1$, so its rank is at least $3$ (in fact exactly $3$)
over every field. The decisive facts (the presented shift, the minimal
state count $2$, the language for all words, and the nonsingular minor) are
formalized and checked in a self-contained Lean~4 project. An independent
Python script confirms them.
\end{abstract}

\section{The conjecture}

The statement (English version) reads:
\begin{quote}
\emph{Definition:} Sofic shifts are factor images of SFTs, presented by
finite graphs with output (hidden Markov models); the Fisher minimal
realization is unique. \emph{Conjecture:} The Fisher minimal state count of
a sofic shift equals the rank of the Hankel matrix of its language; the rank
is computable in polynomial time by a local criterion on words of length
$2n+1$, and the minimal realization graph structure is unique.
\end{quote}
The Chinese version says the same. The conjecture is a conjunction of three
clauses. We refute the first one:
\begin{quote}
(C) \emph{For every sofic shift $X$, the number of states of its Fischer
(minimal right-resolving) presentation equals the rank of the Hankel matrix
of $\Lang(X)$.}
\end{quote}
A conjunction with a false clause is false, so the conjecture is false. The
two later clauses are not needed and we say nothing about them. ``Fisher'' is
a common misspelling of R.~Fischer, who introduced the minimal
right-resolving presentation \cite{fischer}.

\section{Definitions}

We follow Lind and Marcus \cite{lm}. The alphabet is $\{0,1\}$. A point is a
bi-infinite sequence $x=(x_i)_{i\in\mathbb Z}$. A word $w=w_0\cdots w_{n-1}$
\emph{occurs} in $x$ at $i$ if $x_{i+j}=w_j$ for all $j<n$. The
\emph{language} $\Lang(X)$ of a set $X$ of points is the set of words that
occur in some point of $X$.

\paragraph{Labelled graphs and presentations.} A labelled graph $G$ on the
vertex set $\{0,\dots,n-1\}$ is a set of edges $p\xrightarrow{a}q$ with
$a\in\{0,1\}$. The shift $X_G$ presented by $G$ is the set of label sequences
of bi-infinite paths, i.e.\ of all $x$ for which there are vertices $s_i$
($i\in\mathbb Z$) with an edge $s_i\xrightarrow{x_i}s_{i+1}$ for every $i$. A
shift space $X$ is \emph{sofic} if $X=X_G$ for some finite $G$; then $G$
\emph{presents} $X$. (Parallel edges with the same label do not change $X_G$,
so we may take $G$ to be a set of labelled edges.) $G$ is
\emph{right-resolving} if the out-edges of each vertex carry distinct labels.

\paragraph{Fischer cover.} Every irreducible sofic shift has a minimal
right-resolving presentation, unique up to isomorphism \cite{fischer},
\cite[\S3.3]{lm}. It is called the Fischer cover. It has the fewest
vertices among all right-resolving presentations \cite[\S3.3]{lm}. So
the ``Fisher minimal state count'' is
\[
\mathrm{Fis}(X)=\min\{n : X \text{ has a right-resolving presentation with } n \text{ vertices}\}.
\]
We also consider $\mathrm{Pres}(X)=\min\{n : X \text{ has a presentation with } n \text{ vertices}\}$.

\paragraph{Hankel matrix and rank.} The Hankel matrix of a language $\Lang$
is the infinite $0/1$ matrix indexed by pairs of words,
$H[u][v]=1$ if $uv\in\Lang$ and $0$ otherwise. Its row $u$ is the indicator
of the follower set $\F(u)=\{v: uv\in\Lang\}$. The rank of an infinite matrix
over a field $K$ is the dimension of its row space. Equivalently, it is the
supremum of the ranks of its finite submatrices, or the largest $r$ such
that some $r\times r$ minor is nonzero in $K$. This is the meaning of ``Hankel
rank'' in realization theory, from Kronecker's theorem on Hankel matrices of
rational functions to the Fliess and Carlyle--Paz theorems on rational formal
power series and the Hankel rank of hidden Markov processes
\cite{fliess,cp,ito}.

\section{The even shift and its language}

The even shift is
\[
X=\{x\in\{0,1\}^{\mathbb Z} : \text{no block } 0\,1^{2k+1}\,0 \ (k\ge 0) \text{ occurs in } x\}.
\]
A block $0\,1^m\,0$ is exactly a run of $m$ ones bounded by $0$s on both sides,
so $X$ consists of the sequences in which every such run has even length.

\begin{lemma}\label{lem:lang}
For every word $w$: $w\in\Lang(X)$ if and only if $w$ contains no factor
$0\,1^{2k+1}\,0$.
\end{lemma}
\begin{proof}
If $w$ occurs in $x\in X$ at $i$, then a factor $0\,1^{2k+1}\,0$ of $w$ at
position $j$ would occur in $x$ at $i+j$. Conversely, if $w$ has no such
factor, let $x$ be $w$ placed at positions $0,\dots,|w|-1$ and $1$
everywhere else. Every $0$ of $x$ lies inside $w$. So a block
$0\,1^{2k+1}\,0$ of $x$ would lie inside $w$, which is impossible. Hence
$x\in X$ and $w\in\Lang(X)$.
\end{proof}

\section{The Fischer cover has two states}

Let $G$ be the graph with vertices $A=0$, $B=1$ and edges
\[
A\xrightarrow{0}A,\qquad A\xrightarrow{1}B,\qquad B\xrightarrow{1}A.
\]
It is right-resolving and irreducible. It is follower-separated: $0$ labels a
path from $A$ but not from $B$.

\begin{proposition}\label{prop:presents}
$X_G=X$.
\end{proposition}
\begin{proof}
($X_G\subseteq X$.) Let $(s_i)$ be a path with labels $(x_i)$. Suppose
$0\,1^{2k+1}\,0$ occurs at $i$. The only $0$-edge ends at $A$, so
$s_{i+1}=A$. Along the $1$-edges the state alternates, so
$s_{i+1+t}=A$ for even $t$ and $B$ for odd $t$ ($t\le 2k+1$). Hence
$s_{i+2k+2}=B$. But $x_{i+2k+2}=0$ and $B$ has no $0$-edge. This is a contradiction.

($X\subseteq X_G$.) Let $x\in X$. We say position $i$ is in state $B$ if one of
the following holds:
\begin{enumerate}
\item[(a)] there is a $0$ at or after $i$, and the first one is at $i+t$ with $t$ odd;
\item[(b)] there is no $0$ at or after $i$, there is a $0$ before $i$, and the
last one is at $i-1-t$ with $t$ odd;
\item[(c)] $x$ has no $0$ at all and $i$ is odd.
\end{enumerate}
Otherwise $i$ is in state $A$. We check that $s_i\xrightarrow{x_i}s_{i+1}$ is an edge for every $i$.

If $x_i=0$, then (a) fails at $i$ with $t=0$, and (b), (c) fail too. So $s_i=A$.
At $i+1$, (a) would give a block $0\,1^{t}\,0$ at $i$ with $t$ odd, which is
excluded since $x\in X$. (b) fails because the last $0$ before $i+1$ is at
$i$, so $t=0$. (c) fails. So $s_{i+1}=A$ and $A\xrightarrow{0}A$ is an edge.

If $x_i=1$, we show that $s_{i+1}$ is the other state from $s_i$. Then the step
is $A\xrightarrow1B$ or $B\xrightarrow1A$. If there is a $0$ after $i$, say
first at $i+1+t$, then the first $0$ at or after $i$ is at $i+(t+1)$. Only
rule (a) applies at both positions, with parities of $t$ and $t+1$. If there is
no $0$ after $i$ but one before $i$, say last at $i-1-t$, then the last $0$
before $i+1$ is at $(i+1)-1-(t+1)$. Only rule (b) applies, with parities of $t$
and $t+1$. If $x$ has no $0$ at all, rule (c) applies with the parities of $i$
and $i+1$.
\end{proof}

\begin{proposition}\label{prop:lower}
No labelled graph with at most one vertex presents $X$.
\end{proposition}
\begin{proof}
A graph with no vertex has no bi-infinite path, but $0^\infty\in X$. A graph
with one vertex and loop label set $T\subseteq\{0,1\}$ presents the full
shift $T^{\mathbb Z}$. If $0\notin T$, it misses $0^\infty\in X$. If $1\notin T$,
it misses $1^\infty\in X$ (which has no $0$). If $T=\{0,1\}$, it contains
$\cdots000\,1\,000\cdots\notin X$, since this point contains the block $010$.
\end{proof}

\begin{theorem}\label{thm:fis}
$\mathrm{Fis}(X)=\mathrm{Pres}(X)=2$. The graph $G$ is the Fischer cover of $X$.
\end{theorem}
\begin{proof}
By Propositions~\ref{prop:presents} and~\ref{prop:lower}. The
Fischer cover is the unique minimal right-resolving presentation, so it is $G$.
\end{proof}
The verification script also enumerates all $2^8$ labelled graphs on two
vertices. Exactly two of them have the same blocks as $X$ up to length $8$, and
these are $G$ and its relabelling.

\section{The Hankel matrix has rank 3}

By Lemma~\ref{lem:lang}, $H[u][v]=1$ iff $uv$ has no factor $0\,1^{\mathrm{odd}}\,0$.
Take rows $u\in\{11,\,0,\,01\}$ and columns $v\in\{0,\,10,\,1\}$:
\[
\begin{array}{c|ccc}
 & 0 & 10 & 1\\\hline
11 & 110\in\Lang & 1110\in\Lang & 111\in\Lang\\
0 & 00\in\Lang & 010\notin\Lang & 01\in\Lang\\
01 & 010\notin\Lang & 0110\in\Lang & 011\in\Lang
\end{array}
\qquad
M=\begin{pmatrix}1&1&1\\1&0&1\\0&1&1\end{pmatrix}.
\]
We have $\det M=1\cdot(0-1)-1\cdot(1-0)+1\cdot(1-0)=-1$. The integer matrix
$N=\left(\begin{smallmatrix}1&0&-1\\1&-1&0\\-1&1&1\end{smallmatrix}\right)$
satisfies $MN=NM=I$. So $M$ is invertible over every commutative ring, and
in particular over every field.

\begin{theorem}\label{thm:rank}
Over every field $K$, the Hankel matrix of $\Lang(X)$ has rank exactly $3$.
\end{theorem}
\begin{proof}
The rank is at least $3$ because $M$ is a nonsingular submatrix. For the
upper bound we show that $H$ has at most four distinct rows, one of which is
zero. If $u\notin\Lang(X)$, then $\F(u)=\emptyset$, because $\Lang(X)$ is
closed under taking factors. Let $u\in\Lang(X)$. A forbidden block in $uv$ lies
inside $u$ (impossible), inside $v$, or straddles the junction. A straddling
block $0\,1^m\,0$ has its first $0$ in $u$ and all $1$s up to the junction, so
it starts at the last $0$ of $u$. Write $u=u'01^j$ if $u$ contains a $0$, and
$v=1^r0v'$ if $v$ contains a $0$. Then $uv\in\Lang(X)$ iff $v\in\Lang(X)$ and,
when both contain a $0$, $j+r$ is even. Hence $\F(u)=\Lang(X)$ if $u$ has no
$0$, $\F(u)=\F(0)$ if $j$ is even, and $\F(u)=\F(01)$ if $j$ is odd. The rows
lie in the span of the three vectors $\mathbf 1_{\Lang(X)}$, $\mathbf 1_{\F(0)}$
and $\mathbf 1_{\F(01)}$, so the rank is at most $3$.
\end{proof}

\begin{corollary}
$\mathrm{Fis}(X)=2\ne 3=\operatorname{rank}_K H$ for every field $K$. Clause (C), and
with it Conjecture 00000009617, is false.
\end{corollary}

\section{Readings of the clause}

\begin{itemize}
\item \textbf{Rank over a field} ($\mathbb Q,\mathbb R,\mathbb C,\mathbb F_p$, or any other field).
The rank is $3\ne2$. The determinant $-1$ is a unit in every ring, so the
counterexample does not depend on the field. This is the standard meaning of
the rank of a Hankel matrix (Section~2).
\item \textbf{Fisher count.} We read it as the size of the minimal
right-resolving presentation. If it means the least number of states of
\emph{any} presentation, it is still $2$ (Theorem~\ref{thm:fis}). If it means
the left Fischer cover (minimal left-resolving presentation), it is also $2$:
$X$ is closed under reversal, and reversing $G$ gives a left-resolving
presentation; Proposition~\ref{prop:lower} holds for all graphs.
\item \textbf{Indexing of $H$.} Rows and columns may range over all words, over
nonempty words, or over words of length at most $n$ for any $n\ge2$. All these
contain $M$, so the rank is at least $3$ in each case. Indexing by $u^{R}v$ or
transposing does not change the rank.
\item \textbf{Boolean rank (disclosure).} Over the Boolean semiring (OR/AND), the
rank of $H$ is $2$. Indeed $\mathbf 1_{\Lang(X)}=\mathbf 1_{\F(0)}\vee\mathbf 1_{\F(01)}$:
if $v$ has no $0$, then $v\in\F(0)$; otherwise $v\in\F(0)$ or $v\in\F(01)$ by
the parity of its leading run of $1$s. The fooling set $\{(0,0),(01,10)\}$ shows
that the Boolean rank is at least $2$. Under this reading the even shift gives
$2=2$ and is \emph{not} a counterexample. We do not claim anything about the
Boolean reading. It is not the standard meaning: Boolean rank is a different
notion and is always named as such. Over $\mathbb Q$ the vector
$\mathbf 1_{\Lang(X)}$ is not $\mathbf 1_{\F(0)}+\mathbf 1_{\F(01)}$, since both
follower sets contain every word $1^k$. This is why the field rank is $3$.
\item \textbf{Hidden Markov model reading.} One could instead take the Hankel matrix
of the probability function $P(w)$ of a measure on $X$. That matrix is not
determined by the shift or by ``its language''. The clause speaks of the Hankel
matrix of the language, so we use the indicator function of $\Lang(X)$.
\item \textbf{Other covers.} The number of distinct nonempty follower sets of words
of $X$ is $3$. This is the size of the follower-set graph and of the minimal
deterministic automaton of $\Lang(X)$ without its dead state, and it equals
the Hankel rank here. But that graph is not minimal, and it is not the Fischer
cover. The Fischer cover only keeps the follower sets of synchronizing words,
namely $\F(0)$ and $\F(01)$.
\end{itemize}

\section{Lean formalization}
\begingroup\raggedright

The directory \texttt{lean4/} is a self-contained Lean~4.19.0 project. It uses
only the core library (no Mathlib). Points are functions
\texttt{Int}~$\to$~\texttt{Bool} (\texttt{false} is $0$, \texttt{true} is $1$).
\begin{itemize}
\item \texttt{EvenShift x}: no \texttt{ForbiddenAt x i k}, the block
$0\,1^{2k+1}\,0$ at $i$. \texttt{Lang X w}: some point of $X$ contains $w$.
\item \texttt{LGraph n}, \texttt{IsPathLabel}, \texttt{Presents E X}
($X$ equals the set of bi-infinite path labels), \texttt{RightResolving},
\texttt{Sofic}. \texttt{IsFisherCount X m} means that $m$ is the least number
of vertices of a right-resolving presentation. \texttt{IsMinPresCount} is the
same for arbitrary presentations.
\item \texttt{hankelEntry X u v} is $1$ if \texttt{Lang X (u ++ v)} holds and
$0$ otherwise. \texttt{minor}, and \texttt{det} (Laplace expansion).
\texttt{HankelRankAtLeast X r} means that some $r\times r$ minor has nonzero
integer determinant, i.e.\ the rank over $\mathbb Q$ is at least $r$.
\texttt{HankelRankAtLeastMod p} is the same over $\mathbb Z/p$.
\texttt{HankelRankIs X m} means at least $m$ and not at least $m+1$.
\item \texttt{Claim}: for every sofic $X$ and every $m$ with
\texttt{IsFisherCount X m}, \texttt{HankelRankIs X m}. \texttt{ClaimMod p} and
\texttt{ClaimAnyPresentation} are the variants for $\mathbb Z/p$ and for
arbitrary presentations.
\item \texttt{lang\_iff}: Lemma~\ref{lem:lang}, for all words.
\texttt{fisher\_sound}, \texttt{fisher\_complete}, \texttt{fisher\_presents}:
Proposition~\ref{prop:presents}. \texttt{no\_zero\_state}, \texttt{no\_one\_state}:
Proposition~\ref{prop:lower}. \texttt{fisher\_count},
\texttt{min\_presentation\_count}: Theorem~\ref{thm:fis}.
\texttt{fisher\_rightResolving} and \texttt{fisher\_irreducible\_separated}
check the properties of $G$.
\item \texttt{minor\_eq}, \texttt{minor\_det} ($=-1$), \texttt{minor\_unimodular}
(the inverse $N$), \texttt{rows\_independent} (the three rows are linearly
independent over $\mathbb Q$), \texttt{hankel\_rank\_ge\_three} and
\texttt{hankel\_rank\_ge\_three\_mod} (for every $p\ge2$).
\item \texttt{conjecture\_00000009617\_false :}~$\neg$\,\texttt{Claim}. Also
\texttt{conjecture\_00000009617\_false\_mod} ($\forall p\ge 2$, $\neg$\,\texttt{ClaimMod p})
and \texttt{conjecture\_00000009617\_false\_any\_presentation}.
\texttt{fisher\_lt\_rank}: every $m$ with \texttt{IsFisherCount X m} equals $2$, and the rank is
at least $m+1$.
\end{itemize}
Non-vacuity: \texttt{evenShift\_sofic} and \texttt{fisher\_count} show that the
hypotheses of \texttt{Claim} hold for $X$. \texttt{allZero\_mem},
\texttt{allOne\_mem} and \texttt{single\_one\_not\_mem} show that $X$ is neither
empty nor full. The final examples show that \texttt{Lang} is nontrivial and
that \texttt{det} returns $0$ on a singular matrix.

\paragraph{Limitations.} Lean proves rank \emph{at least} $3$, which suffices.
The upper bound in Theorem~\ref{thm:rank} and the Boolean-rank remark are proved
in this report and checked on finite truncations by \texttt{verify.py}, but not
in Lean. Lean works over the alphabet $\{0,1\}$. The clause quantifies over all
sofic shifts, so a binary counterexample suffices. The project has no
\texttt{sorry}, no \texttt{native\_decide} and no added axioms.
\texttt{\#print axioms} reports only the standard \texttt{propext},
\texttt{Classical.choice} and \texttt{Quot.sound}. Classical logic is used to
define the path in \texttt{fisher\_complete} and the Hankel entries of an
arbitrary shift.

\endgroup

\section{Independent check}

\texttt{verify.py} (Python~3 standard library) does the following:
\begin{itemize}
\item It checks that three descriptions of $\Lang(X)$ agree on all words of length
at most $14$: forbidden blocks, parities of inner runs, and paths in $G$.
\item It enumerates all labelled graphs with $0$, $1$ or $2$ vertices. None with
fewer than $2$ vertices matches, and the matching $2$-vertex graphs are $G$ up
to relabelling.
\item It recomputes the minor, its determinant $-1$ and its inverse.
\item It computes the rank of the Hankel truncation on all $127$ words of length at
most $6$: it is $3$ over $\mathbb Q$ and over $\mathbb F_2,\mathbb F_3,\mathbb F_5,\mathbb F_7$.
\item It confirms that only $4$ distinct rows occur.
\item It confirms the Boolean-rank identity.
\end{itemize}

\section{Reproduction}
\begin{verbatim}
cd lean4 && lake build       # Lean 4.19.0, core only
python3 verify.py            # standard library only
pdflatex report.tex
\end{verbatim}

\begin{thebibliography}{9}
\bibitem{fischer} R.~Fischer, Sofic systems and graphs, \emph{Monatsh. Math.}
80 (1975), 179--186.
\bibitem{lm} D.~Lind and B.~Marcus, \emph{An Introduction to Symbolic Dynamics
and Coding}, Cambridge University Press, 1995 (2nd ed.\ 2021).
\bibitem{fliess} M.~Fliess, Matrices de Hankel, \emph{J. Math. Pures Appl.}
53 (1974), 197--224.
\bibitem{cp} J.~W.~Carlyle and A.~Paz, Realizations by stochastic finite
automata, \emph{J. Comput. System Sci.} 5 (1971), 26--40.
\bibitem{ito} H.~Ito, S.-I.~Amari and K.~Kobayashi, Identifiability of hidden
Markov information sources and their minimum degrees of freedom,
\emph{IEEE Trans. Inform. Theory} 38 (1992), 324--333.
\end{thebibliography}

\end{document}
